% Groente- en fruitkalender
% Starts on a left/verso page (not \cleardoublepage's usual right/recto) so
% the Groente table below and the Fruit table after it (page break via
% \clearpage further down) land on facing pages of the same spread.
\cleartoevenpage
% Seizoenskalender: a \subchapter-level (H2) heading grouped under the
% single "Inleiding" chapter-level heading in main.tex.
\subchapter{Seizoenskalender}
{\itshape\small\color{inkfaint}Wanneer onze groente en al het fruit op hun best zijn.\par}\vspace{1.2em}

\lettrine{G}{roente} kopen in het seizoen is meestal lekkerder, en ook nog
goedkoper. Hieronder staan alleen de groenten die we in dit boek ook echt
gebruiken, anders wordt het een lijstje waar je toch niets aan hebt. Fruit
eten we het hele jaar door gewoon los, dus daarvan staat de volledige lijst
erbij.

% Local, appendix-only helpers: draw one repeatable timeline row per
% groente/fruit instead of hand-typing a description for each. \seizoenbalk
% draws the J-D month grid with the in-season stretch(es) filled in;
% \seizoenrij wraps that into one table row so adding a new item is a
% single line, with no manual row/column bookkeeping. \seizoenmaanden draws
% the matching J F M A M J J A S O N D header, repeated automatically by
% longtable on every page the table spans.
\newcommand{\seizoenmaanden}{%
  \begin{tikzpicture}[x=0.75cm]
    % Pin the bounding box to the same 0--12 span as \seizoenbalk's grid;
    % without this, the box is fit to the node text instead, whose default
    % center anchor shifts the whole row left of the grid it's meant to
    % label.
    \useasboundingbox (0,-0.3) rectangle (12,0.3);
    \foreach \m [count=\i from 0] in {J,F,M,A,M,J,J,A,S,O,N,D} {
      \node[font=\tiny\itshape, text=inkfaint] at (\i+0.5,0) {\m};
    }
  \end{tikzpicture}}
\newcommand{\seizoenbalk}[1]{%
  \begin{tikzpicture}[x=0.75cm, y=0.045cm]
    \draw[line width=0.3pt, color=inkfaint] (0,0) -- (12,0);
    \foreach \i in {0,...,12} {
      \draw[line width=0.3pt, color=inkfaint] (\i,-1) -- (\i,1);
    }
    \foreach \start/\eind in {#1} {
      \draw[line width=3.4pt, color=terracotta, line cap=round]
        (\start-1,0) -- (\eind,0);
    }
  \end{tikzpicture}}
\newcommand{\seizoenrij}[2]{#1 & \seizoenbalk{#2} \\}
\newcommand{\seizoenkop}{%
  & \seizoenmaanden \\
  \multicolumn{2}{@{}l@{}}{\textcolor{terracotta}{\rule{\linewidth}{0.4pt}}} \\[0.1em]}

\blockrule{Groente}
{\footnotesize\renewcommand{\arraystretch}{1.05}
\begin{longtable}{r@{\hspace{0.45em}}l}
\seizoenkop
\endhead
\seizoenrij{Andijvie}{4/12}
\seizoenrij{Artisjok}{6/10}
\seizoenrij{Asperge}{4/6}
\seizoenrij{Aubergine}{7/10}
\seizoenrij{Bleekselderij}{7/11}
\seizoenrij{Bloemkool}{3/11}
\seizoenrij{Boerenkool}{1/3,10/12}
\seizoenrij{Broccoli}{6/11}
\seizoenrij{Courgette}{7/11}
\seizoenrij{Doperwten}{5/7}
\seizoenrij{Komkommer}{7/10}
\seizoenrij{Maïs}{8/9}
\seizoenrij{Paddenstoelen}{1/12}
\seizoenrij{Paprika}{7/10}
\seizoenrij{Pastinaak}{1/3,10/12}
\seizoenrij{Peultjes}{5/7}
\seizoenrij{Pompoen}{1/4,8/12}
\seizoenrij{Prei}{1/12}
\seizoenrij{Rodekool}{1/3,7/12}
\seizoenrij{Sperzieboon}{6/10}
\seizoenrij{Spinazie}{3/10}
\seizoenrij{Spitskool}{5/10}
\seizoenrij{Spruiten}{1/3,10/12}
\seizoenrij{Tomaat}{6/10}
\seizoenrij{Ui}{1/12}
\seizoenrij{Venkel}{6/11}
\seizoenrij{Witlof}{1/4,10/12}
\seizoenrij{Wittekool}{1/3,7/12}
\seizoenrij{Wortel}{1/12}
\end{longtable}}

% Nog niet in dit boek gebruikt, dus hier alleen als commentaar. Zodra een
% nieuw recept (/add-recipe) zo'n groente gebruikt, haal de bijbehorende
% regel hierboven uit commentaar en zet 'm op alfabet tussen de actieve
% \seizoenrij-regels in de Groente-tabel, en verwijder 'm hieronder. De
% regel staat hieronder al in tabelvorm, dus uitcommentariëren + verplaatsen
% is genoeg.
% Seizoenen zijn uit dezelfde Voedingscentrum-kalender als hierboven.
% \seizoenrij{Aardpeer}{1/1,11/12}
% \seizoenrij{Chinese kool}{7/11}
% \seizoenrij{Groene selderij}{1/12}
% \seizoenrij{Knolselderij}{1/3,8/12}
% \seizoenrij{Koolraap}{1/1,11/12}
% \seizoenrij{Koolrabi}{5/10}
% \seizoenrij{Kropsla}{5/10}
% \seizoenrij{Meiraap}{1/6}
% \seizoenrij{Paksoi}{4/10}
% \seizoenrij{Raap}{9/12}
% \seizoenrij{Raapsteel}{3/5,8/10}
% \seizoenrij{Rabarber}{3/7}
% \seizoenrij{Radijs}{3/10}
% \seizoenrij{Rammenas}{1/3,9/12}
% \seizoenrij{Rode biet}{1/12}
% \seizoenrij{Roodlof}{3/5,10/10}
% \seizoenrij{Savooikool}{1/3,6/12}
% \seizoenrij{Schorseneer}{1/3,9/12}
% \seizoenrij{Snijbiet}{3/11}
% \seizoenrij{Snijboon}{7/10}
% \seizoenrij{Tuinboon}{6/7}
% \seizoenrij{Veldsla}{1/3,10/12}
% \seizoenrij{Waterkers}{1/1,4/12}
% \seizoenrij{Winterpostelein}{1/3,10/12}

\clearpage
\blockrule{Fruit}
{\footnotesize\renewcommand{\arraystretch}{1.05}
\begin{longtable}{r@{\hspace{0.45em}}l}
\seizoenkop
\endhead
\seizoenrij{Aardbei}{4/8}
\seizoenrij{Abrikoos}{6/8}
\seizoenrij{Appel}{1/5,8/12}
\seizoenrij{Blauwe bes}{6/9}
\seizoenrij{Braam}{7/10}
\seizoenrij{Druif}{8/11}
\seizoenrij{Framboos}{6/10}
\seizoenrij{Kers}{6/8}
\seizoenrij{Kiwibes}{9/10}
\seizoenrij{Kweepeer}{9/10}
\seizoenrij{Meloen}{8/11}
\seizoenrij{Nectarine}{6/9}
\seizoenrij{Peer}{1/5,8/12}
\seizoenrij{Perzik}{6/9}
\seizoenrij{Pruim}{8/9}
\seizoenrij{Rode bes}{6/7}
\seizoenrij{Vijg}{8/9}
\seizoenrij{Zwarte bes}{6/8}
\end{longtable}}

\vspace{0.8em}
{\footnotesize\itshape\color{terracotta}Bron: Voedingscentrum, \emph{Seizoengroente- en
fruitkalender}\cite{voedingscentrum-seizoenskalender}.\par}
